\chapter{Implementation}
\label{ch:implementation}

\section{System architecture}
\label{sec:architecture}

The study is a piece of software before it is a result. The software had to put
four different language models, three reached over the network and one running on
a local GPU, into the same repeated game, record every message and move, and stop
before a hundred euros were spent. The nearest related framework,
FAIRGAME~\cite{fairgame2025}, is built for two-player, single-model, matrix-form
games, and none of those three assumptions holds here, so the system was written
from scratch. Three of FAIRGAME's patterns were worth keeping and were
reimplemented rather than inherited: an experiment specified by a JSON
configuration and a seed, prompts built from text templates, and a retry loop that
tolerates malformed model output.

The code is a small Python package with a deliberate dependency direction. A
configuration layer turns a JSON experiment description and an integer seed into
typed objects. Above it sit two independent abstractions. A \emph{game} owns the rules: it renders each agent's prompt, checks
the action that comes back, and supplies a fallback for when none is valid. An
\emph{agent} owns a single model and the act of calling it, behind one interface
that all four backends share. A runner drives the two together round by round, and
beneath them a budget tracker and a SQLite log watch the spending and record every
call.

The architecture turns on two choices. The first is that the four
backends share one pipeline: retrying, logging, and the budget check live once in
the agent base class, and a backend implements only the raw call to its provider,
so the retry and cost logic is written and tested once rather than four times. The
second is the separation of two kinds of failure. A transport error, a rate limit
or a timeout, is retried with backoff and, if it persists, raised rather than
swallowed, because a dead connection is not a model giving a bad answer and
treating it as one would corrupt the data. A malformed or invalid response is the
other kind: the agent re-prompts a few times and, failing that, takes the game's
safe default action (Section~\ref{sec:protocols}). The two are counted separately,
so only genuine parse failures bear on the data-quality gate.

One more choice follows from the hardware. The local model answers one request at a
time on a single consumer GPU, so a round gains nothing from parallel local calls;
instead the agents' calls are fanned out across threads, which lets the paid
network calls overlap the serial local queue, and the results are gathered back in
their original order so the concurrency never changes the outcome. The log holds it
together. Every call is written to SQLite the moment it returns, with
its prompt, response, parsed action, cost, and failure counts; the budget tracker
reads the running total back from that log to enforce the cap, and the metrics
later read it to reconstruct what each agent said against what it did.

\section{Agent abstraction and model backends}
\label{sec:agents}

Every model in the study is reached through one small interface, so that the
calling code never depends on which provider answered. An agent's entry point,
\texttt{act}, takes a rendered prompt, a validator, and a fallback action, and
returns a validated action with the accounting the log needs; a backend supplies
only the raw call. Listing~\ref{lst:agent} is the pipeline they share.

\begin{lstlisting}[caption={The shared agent pipeline from \texttt{agents/base.py} (docstrings abridged): a backend implements only the raw model call, while \texttt{act} owns retrying, validation, fallback, and logging.}, label={lst:agent}]
class BaseAgent(ABC):
    @abstractmethod
    def _complete(self, prompt: str, *, json_mode: bool) -> CompletionResult:
        """Make one raw model call and return text plus cost accounting."""

    def act(self, prompt: str, validator: Validator, fallback: Action, *,
            json_mode: bool = True, experiment_id: str | None = None,
            round_no: int | None = None) -> ActionResult:
        result = ActionResult(action=fallback)
        last_raw: str | None = None
        for attempt in range(self.max_parse_retries + 1):
            result._attempts = attempt + 1
            completion = self._call_with_backoff(prompt, json_mode=json_mode, result=result)
            last_raw = completion.raw_text
            result.input_tokens += completion.input_tokens
            result.output_tokens += completion.output_tokens
            result.cost_eur += completion.cost_eur
            try:
                payload = _parse_json_object(completion.raw_text)
                action = validator(payload)
            except (json.JSONDecodeError, ValueError, KeyError, TypeError):
                result.parse_fail_count += 1
                continue
            result.action = action
            result.message = payload.get("message") if isinstance(payload, dict) else None
            result.parsed_action = action
            result.parse_ok = True
            break
        else:
            # Loop exhausted without break: every attempt failed to parse.
            result.fallback_used = True
            result.parsed_action = fallback
        result.raw_response = last_raw
        self._log(result, prompt, experiment_id, round_no)
        return result
\end{lstlisting}
\source{Author}

Each pass calls the backend through a wrapper that checks the budget and retries
transport failures with backoff; if the API stays down the wrapper re-raises, so
an exhausted transport failure propagates out of \texttt{act} rather than being
logged as a move. When a response comes back it is parsed and handed to the game's
validator. A malformed or rejected response increments the parse counter and the
loop tries again, up to a fixed number of attempts; if none succeeds, the
\texttt{else} branch sets the fallback flag and the call resolves to the game's
safe default action (Section~\ref{sec:protocols}). That outcome, parsed or fallen
back, is logged before the result is returned, and the single abstract method is
the only thing a backend provides.

The four backends differ only in that call. Llama runs locally through Ollama,
which is asked for JSON at the prompt level and incurs no API cost because it runs
on the machine's own GPU. The three hosted models cost real money, converted to
euros and recorded per call; the Claude backend also prefills the assistant's
reply with an opening brace, so the model emits a JSON object rather than prose,
which fixed a parse failure seen in the first paid pilot.

\section{The game engine}
\label{sec:games}

A game owns the rules and nothing else. It exposes one small interface that the
runner drives without knowing which game it holds: \texttt{reset} starts a fresh
match, \texttt{render\_prompt} turns the current state and the history into the text
an agent reads, \texttt{validator} returns a function that checks one parsed response
and extracts the move, \texttt{fallback\_action} supplies the safe default for a
response that never parses, and \texttt{step} resolves a round into payoffs and
advances the state. The Public Goods Game and the Trust Game each implement that
interface, and the agent layer of Section~\ref{sec:agents} never refers to either by
name.

The round cycle is the same for both. The game renders a prompt, the agent returns
JSON, the validator extracts the one field that is a legal move, and \texttt{step}
turns the round's moves into payoffs. Listing~\ref{lst:pgg} is the Public Goods half
of that cycle, with the validator, the fallback, and the payoff rule together.

\begin{lstlisting}[caption={The Public Goods decision cycle (\texttt{games/public\_goods.py}, the \texttt{step} docstring abridged): the validator accepts only an integer contribution in range and returns it alone; the fallback free-rides; \texttt{step} multiplies the pool, splits it, and pays each agent its share plus what it kept.}, label={lst:pgg}]
    def validator(self, agent_id: str, state: dict[str, Any]) -> Validator:
        """Return a validator accepting ``contribution`` in ``[0, endowment]``."""
        endowment = self.endowment

        def _validate(payload: dict[str, Any]) -> Action:
            if not isinstance(payload, dict):
                raise ValueError("expected a JSON object")
            contribution = payload["contribution"]
            # bool is a subclass of int; reject it explicitly.
            if isinstance(contribution, bool) or not isinstance(contribution, int):
                raise ValueError("contribution must be an integer")
            if not 0 <= contribution <= endowment:
                raise ValueError(f"contribution out of range [0, {endowment}]")
            return {"contribution": contribution}

        return _validate

    def fallback_action(self, agent_id: str, state: dict[str, Any]) -> Action:
        """Safe default when an agent never parses: contribute 0 (free-ride)."""
        return {"contribution": 0}

    def step(
        self, actions: dict[str, Action]
    ) -> tuple[dict[str, Any], dict[str, float], bool]:
        """Resolve one round and advance state."""
        total = sum(action["contribution"] for action in actions.values())
        pool = total * self.multiplier
        share = pool / self.n
        payoffs = {
            agent_id: share + (self.endowment - action["contribution"])
            for agent_id, action in actions.items()
        }
        done = self._round >= self.rounds
        self._round += 1
        return self._public_state(), payoffs, done
\end{lstlisting}
\source{Author}

The validator is what keeps the rest of the system honest. It reads the one field it
cares about, rejects a boolean, which Python would otherwise accept as an integer, and
any out-of-range value, and returns a fresh \texttt{\{"contribution": n\}}. Anything
else the model emitted, including a message in the communication condition, is dropped
here and never enters the game state. That is the first of the two reasons an agent
never sees another's words: the move that gets stored is only the number. The second is
the history the next prompt is built from, which carries each past contribution and no
message at all (Section~\ref{sec:communication}). When a response cannot be parsed even
after the agent's retries, \texttt{fallback\_action} stands in with a contribution of
zero, the free-riding default, so a broken call costs the agent rather than corrupting
the round.

\texttt{step} closes the round. It sums the contributions, multiplies the pool by the
factor \(1.6\), splits it evenly, and pays each agent that share plus the endowment it
kept; the game ends once the tenth round resolves. The numbers behind all of this, four
players, an endowment of twenty, the multiplier, ten rounds, are read from the JSON
config (Section~\ref{sec:protocols}), not built into the engine.

The Trust Game implements the same interface over a different rule. It is sequential, so
it is mapped onto the simultaneous runner by splitting each of the ten rounds into two
turns, an investor's send and a trustee's return, with the idle side taking a no-op; the
investor ends with the endowment less what it sent plus what came back, and the trustee
keeps three times the send less what it returned. Which tier opens as investor follows
the order of agents in the config, which is how a composition places a capability tier on
a side of the exchange. As in the Public Goods Game, the history a later prompt shows
carries only the amounts sent and returned, never a message.

\section{The communication condition (cheap talk)}
\label{sec:communication}

Communication is a property of the prompt, not a separate channel between players.
When it is on, the game's instruction asks each agent to write a short public
message along with its action and tells the agent the others will see it. The
message is real text the agent produces, and it is captured and stored. But the
runner does not pass it on: each round, only the agents' actions enter the history
that builds the next round's prompts, as Listing~\ref{lst:round-record} shows.

\begin{lstlisting}[caption={Recording a round (\texttt{runner/tournament.py}): each message is written to the database alongside the action, but the \texttt{history} that renders the next prompt holds only actions.}, label={lst:round-record}]
actions = {aid: result.action for aid, result in result_by_id.items()}
for aid in agent_ids:
    result = result_by_id[aid]
    total_cost += result.cost_eur
    db.insert_action(config.db_path, db.ActionRecord(
        experiment_id=experiment_id, round_num=round_num, agent_id=aid,
        action_json=result.action, message=result.message,
        retries=result.parse_fail_count + result.api_error_count,
        cost_eur=result.cost_eur, fallback_used=result.fallback_used))
history.append({"round": round_num, "actions": actions})
\end{lstlisting}
\source{Author}

The message and the history go to different places. Each message is written to the
database on the action's row, where the deception measures and the message
classifier later read it (Sections~\ref{subsec:deception}
and~\ref{subsec:classification}). The \texttt{history} list, by contrast,
accumulates only the round's actions, and it is that list the game renders into
the next prompt. No step copies a message into the history or into any agent's
prompt. An agent therefore sees what the others did in past rounds, never what
they said.

This makes the condition a manipulation of framing rather than of exchange. Both
conditions already show an agent the others' past contributions or transfers; the
one thing communication adds is the request to make a public pledge, together with
the prompt's claim that it will be read. Whatever the condition changes in
behavior, it changes through that framing, since no pledge is ever delivered to
anyone. The prompt's wording overstates what the runner provides, and the results
chapter reads the effect on that understanding (Chapter~\ref{ch:results}). The
messages keep their other role intact: the deception measures of
Section~\ref{subsec:deception} compare an agent's own message with its own action,
and so do not depend on another player receiving it.

\section{The metric-computation pipeline}
\label{sec:metric-pipeline}

The metrics never call the players. They read the SQLite log the runner wrote and
reconstruct, for each agent, what it said against what it did. Most are pure functions
over the logged numbers: the Gini coefficient and the Atkinson index run on the payoff
vectors, the capability premium subtracts two group means, and a ten-thousand-resample
bootstrap puts a confidence interval on the Gini. The deception measure is the one that
needs more than arithmetic, because turning a free-text message into something
comparable to an action takes a second model.

For every round in which an agent sent a message, a separate Haiku call reads that
message, together with a neutral description of the game and never the agent's actual
move, and returns a \emph{stated policy}: a probability distribution over the game's
discrete action support, the five binned contributions \(\{0,5,10,15,20\}\) in the
Public Goods Game and the six investor sends in Trust. Listing~\ref{lst:stated} is that
call.

\begin{lstlisting}[caption={Eliciting the stated policy (\texttt{metrics/deception.py}, docstring abridged): one Haiku call per message turns the text into a distribution \(q\) over the action support, validated and renormalized, with a uniform fallback on parse failure.}, label={lst:stated}]
def extract_stated(
    message: str,
    game_context: str,
    support: Sequence[int],
    *,
    agent: AnthropicAgent,
) -> StatedResult:
    """Extract a stated-policy distribution from one message via Haiku."""
    support = tuple(support)
    prompt = _SYSTEM_TEXT + "\n\n" + _USER.render(
        context=game_context or "(none)",
        support=", ".join(str(v) for v in support),
        message=message,
    )
    result = llm_score(
        agent,
        prompt,
        validator=distribution_validator(support),
        fallback=uniform_fallback(support),
    )
    return StatedResult(
        distribution=dict(result.action["distribution"]),
        cost_eur=result.cost_eur,
        fallback_used=result.fallback_used,
        raw_response=result.raw_response,
    )
\end{lstlisting}
\source{Author}

The scorer is deliberately blind to behavior: it is handed the message and the game,
never the action, so the policy it returns depends on the words alone. From that
distribution \(q\) and the agent's own moves, the pipeline computes two scores,
Listing~\ref{lst:surprisal}.

\begin{lstlisting}[caption={The two deception scores (\texttt{metrics/deception.py}, the core of \texttt{deception\_scores}; the per-round export detail abridged): for each round with a message, the headline policy-KL and the surprisal \(-\log q(a_t)\) of the action actually taken, each averaged over the scored rounds.}, label={lst:surprisal}]
    support = tuple(support)
    revealed_global = empirical_distribution(revealed_actions, support)
    revealed_smoothed = smooth(revealed_global, epsilon)

    kls: list[float] = []
    surprisals: list[float] = []
    for action, stated in zip(revealed_actions, stated_per_round):
        if stated is None:
            continue
        stated_arr = np.asarray([stated[value] for value in support], dtype=float)
        stated_smoothed = smooth(stated_arr, epsilon)
        if direction == "stated_given_revealed":
            kls.append(kl_divergence(stated_arr, revealed_smoothed))
        else:
            kls.append(kl_divergence(revealed_global, stated_smoothed))
        idx = _action_index(action, support)
        surprisals.append(-math.log(stated_smoothed[idx]))

    n = len(kls)
    return {
        "kl_headline": float(np.mean(kls)) if n else float("nan"),
        "surprisal": float(np.mean(surprisals)) if n else float("nan"),
        "n_rounds": n,
        "direction": direction,
        # ... per-scored-round (kl, surprisal) detail omitted
    }
\end{lstlisting}
\source{Author}

The two scores read the same words differently. The \emph{surprisal} is \(-\log
q(a_t)\), the negative log probability the stated policy gave to the action the agent
actually took that round, averaged over the rounds that carried a message. It is
contemporaneous: round \(t\)'s words are tested against round \(t\)'s move, so a high
value means an agent did something its own message had made unlikely. The
\emph{policy-KL} compares distributions instead, \(D_{\mathrm{KL}}(p_{\mathrm{revealed}}
\,\|\, q)\), but \(p_{\mathrm{revealed}}\) is the agent's binned action distribution over
the \emph{whole} game rather than that round. Pooling that way rewards a peaked
distribution: an agent that reliably plays the same cooperative move has a sharp
\(p_{\mathrm{revealed}}\) that diverges hard from any hedged stated distribution, so the
steadiest cooperators score as the most deceptive, which is backwards. Because it
measures how concentrated a policy is rather than whether words match deeds, the
analysis keeps the policy-KL only as a negative control and treats the surprisal as the
primary signal (Section~\ref{subsec:deception}).

Two things about this code are worth stating plainly. Its own names lag the analysis:
the stored field is \texttt{kl\_headline} and the surprisal is labelled a companion,
wording that predates the pre-registration's choice of the surprisal as the primary,
exploratory measure and the policy-KL as the negative control; the thesis follows the
pre-registration, not the field name. And the stated policy comes from Haiku, the same model family that also plays as a
strong agent, so the scorer is in part one of the players; that coupling and its
rebuttal are set out in Section~\ref{sec:integrity}. Both scores smooth the stated
distribution first (\(\varepsilon = 10^{-3}\)) so a zero can never send the logarithm to
infinity, and both are in nats, the logs being natural. Because each compares an agent's
own message with its own action, neither depends on a message reaching another player,
and so neither is touched by the fact that messages are not delivered
(Section~\ref{sec:communication}).

\section{Reproducibility: seeding, isolation, resumability, cost logging}
\label{sec:repro}

A run that takes days under a hard budget cannot be the kind of thing that starts over
when a laptop sleeps or a provider times out. Four properties keep it restartable and
reproducible, and each is mechanical rather than a promise.

\emph{Seeding.} An experiment is fully specified by its JSON config and one integer
seed, both written into the database when the game starts. The seed drives the single
place randomness enters a game: the rendered history shuffles the other players' past
contributions each time it is built, so their order on the page carries no identity an
agent could track from one round to the next. The scope is worth being exact about. The
seed fixes that anonymizing shuffle and is stored with the run; it does not make a
language model's text reproducible, which no seed can. The same config and seed give the
same game setup, the same anonymization, and the same record to analyze, not the same
sentences from a model.

\emph{Isolation and resumability.} Every game writes under its own experiment id, a
random hex string that is part of the primary key of every row it produces, so games
that share one database never collide and any one of them can be deleted without
disturbing the rest. That is what makes a clean resume possible. A finished game has an
end time on its row; a game whose process died mid-run has none. Listing~\ref{lst:resume}
is the resume filter: it deletes every unfinished game across all the tables so it can
run again without leaving duplicate rows, and it skips every cell already finished, keyed
on the config name that encodes the cell and its seed.

\begin{lstlisting}[caption={The idempotent resume (\texttt{scripts/run\_production.py}, a status print abridged): partial runs are purged row by row, finished cells are recognized by config name and skipped, so re-running the matrix completes only what is missing.}, label={lst:resume}]
def _resume_filter(configs: list[ExperimentConfig], db_path: str) -> list[ExperimentConfig]:
    """For --resume: purge partial (unfinished) experiments and drop already-finished cells.

    A finished experiment has ``end_time`` set; its config ``name`` (cell+seed) is the
    idempotency key. Partial games (process died mid-run) are deleted across all tables
    so they re-run cleanly with no duplicate rows.
    """
    import json

    con = db.connect(db_path)
    try:
        partial = [r[0] for r in con.execute(
            "SELECT id FROM experiments WHERE end_time IS NULL").fetchall()]
        for eid in partial:
            for table in ("agents", "rounds", "actions", "payoffs", "llm_calls"):
                con.execute(f"DELETE FROM {table} WHERE experiment_id = ?", (eid,))
            con.execute("DELETE FROM experiments WHERE id = ?", (eid,))
        con.commit()
        finished = {
            json.loads(r[0])["name"]
            for r in con.execute(
                "SELECT config_json FROM experiments WHERE end_time IS NOT NULL").fetchall()
        }
    finally:
        con.close()
    remaining = [c for c in configs if c.name not in finished]
    # ... log how many were kept, purged, and remain
    return remaining
\end{lstlisting}
\source{Author}

A note the listing earns: this isolation is by id inside one shared database, not a
separate file per game, and the resume lives in the production script, not in the pilot
runner, which has no skip logic and always replays from its base seed.

\emph{Cost logging and the cap.} Every call an agent makes is written to the log the
moment it returns, with its cost in euros, and the budget tracker reads the running total
straight back from that log when it starts, so the cap is cumulative across every game in
the database and survives a restart. Before each call the tracker compares the total with
the limit and raises rather than spend past it, so the call that would cross the line is
never made; the tracker introduced in Section~\ref{sec:architecture} is, here, the thing
that lets a run halt and resume without losing count. The ceiling written into the code is
one hundred euros; the production run was given a tripwire of eighty-five, set as a margin
below the ceiling, and the final spend is reported with the results
(Chapter~\ref{ch:results}).
